\item \textbf{ACTIVIDAD: El comportamiento del caos en sistemas armónicos.}
Divida a su grupo en dos mitades. Cada subgrupo debería trabajar en una de las siguientes situaciones:
\begin{itemize}
    \item \textbf{Situación (i):} Un resorte colgante se estira en $35.0\text{ cm}$ cuando un objeto de masa $450\text{ g}$ se cuelga en reposo. En esta situación, definimos su posición como $x = 0$, con $x$ positiva hacia arriba. El objeto se baja $18\text{ cm}$ adicionales y se libera del reposo para oscilar sin fricción.
    \item \textbf{Situación (ii):} Otro resorte colgante se estira $35.5\text{ cm}$ cuando un objeto de masa $440\text{ g}$ se cuelga en reposo. Definimos esta nueva posición como $x = 0$. Este objeto se baja $18\text{ cm}$ adicionales y se libera del reposo para oscilar sin fricción.
\end{itemize}
\begin{enumerate}
    \item Para cada una de estas situaciones, responda las dos preguntas siguientes:
    \begin{enumerate}
        \item ¿Cuál es la posición $x$ del objeto en un momento $84.4\text{ s}$ después?
        \item ¿Qué distancia total ha recorrido el objeto vibrante en el intervalo de tiempo de $84.4\text{ s}$?
    \end{enumerate}
    Cuando finalicen los cálculos, compare los resultados de las dos situaciones.
    \item ¿Por qué las respuestas a la pregunta (a) son tan diferentes cuando los datos iniciales en las situaciones (i) y (ii) son muy similares, mientras que las respuestas a la pregunta (b) son relativamente similares?
    \item ¿Esta circunstancia revela una dificultad fundamental para calcular el futuro?
\end{enumerate}